\documentclass[12pt, letterpaper]{article}

% --- Paquetes de Idioma y Codificación ---
\usepackage[utf8]{inputenc}
\usepackage[T1]{fontenc}
\usepackage[spanish, es-tabla]{babel}

% --- Paquetes de Matemáticas y Física ---
\usepackage{amsmath}
\usepackage{amssymb}
\usepackage{physics} % Proporciona comandos útiles para notación de Dirac, vectores, etc.
\usepackage[version=4]{mhchem} % Para fórmulas químicas elegantes (\ce{Ti3C2O2})

% --- Configuración de Márgenes (Estándar Académico) ---
\usepackage{geometry}
\geometry{
    letterpaper,
    top=2.5cm,
    bottom=2.5cm,
    left=3.0cm,
    right=2.5cm
}

% --- Paquetes de Tablas y Figuras ---
\usepackage{booktabs} % Para líneas horizontales elegantes en tablas (\toprule, \midrule, \bottomrule)
\usepackage{graphicx}
\usepackage{float}     % Permite fijar la posición de figuras/tablas con [H]
\usepackage{caption}   % Personalización de pies de foto y títulos de tablas
\captionsetup{font=small, labelfont=bf}

% --- Configuración de Ruta para Imágenes ---
\graphicspath{{figuras/}}

% --- Bibliografía (BibLaTeX) ---
\usepackage{csquotes}  % Recomendado por biblatex al usar babel en español
\usepackage[
    backend=biber,
    style=numeric-comp, % Estilo numérico compacto comúnmente usado en física
    sorting=none,       % Orden de aparición
    doi=true,
    isbn=false,
    url=false
]{biblatex}
\addbibresource{referencias.bib}

% --- Enlaces y Metadatos del PDF ---
\usepackage{hyperref}
\hypersetup{
    colorlinks=true,
    linkcolor=blue,
    citecolor=blue,
    filecolor=magenta,      
    urlcolor=cyan,
    pdftitle={Ingeniería de interfaces y modulación topológica en heteroestructuras Heusler/MXene},
    pdfauthor={Saúl Ángel Sánchez González},
    pdfsubject={Protocolo de Investigación Doctoral},
    pdfkeywords={Heusler, MXene, DFT, interfaces, física topológica}
}

% --- Estilo de Página ---
\usepackage{fancyhdr}
\pagestyle{fancy}
\setlength{\headheight}{15pt} % Ajustado para evitar advertencias de fancyhdr
\fancyhf{}
\fancyhead[L]{\small\itshape Ingeniería de interfaces en Heusler/MXene}
\fancyhead[R]{\small\thepage}
\renewcommand{\headrulewidth}{0.4pt}
\renewcommand{\footrulewidth}{0pt}

% --- Cuerpo del Documento ---
\begin{document}

% ==================== PORTADA ====================
\begin{titlepage}
    \centering
    \vspace*{0.5cm}
    
    {\scshape\LARGE Universidad Autónoma de Nuevo León} \\
    \vspace{0.4cm}
    {\large Facultad de Ciencias Físico Matemáticas} \\
    \vspace{0.2cm}
    {\large Subdirección de Estudios de Posgrado} \\
    
    \vspace{1.5cm}
    
    % Línea decorativa
    \rule{\linewidth}{1mm} \\
    \vspace{0.4cm}
    {\Large\bfseries INGENIERÍA DE INTERFACES Y MODULACIÓN TOPOLÓGICA EN HETEROESTRUCTURAS HEUSLER/MXENE} \\
    \vspace{0.3cm}
    \rule{\linewidth}{1mm} \\
    
    \vspace{1.5cm}
    
    {\large\bfseries Protocolo de Investigación Doctoral} \\
    \vspace{0.3cm}
    {\large Presentado para cumplir parcialmente con los requisitos del programa de:} \\
    {\large\itshape Doctorado en Ingeniería Física} \\
    
    \vspace{2.0cm}
    
    \begin{minipage}{0.85\textwidth}
        \centering
        \large
        \textbf{Autor:} \\
        Saúl Ángel Sánchez González \\
        \vspace{0.8cm}
        \textbf{Director de Tesis:} \\
        Dr. Carlos Luna
    \end{minipage}
    
    \vfill
    
    \rule{0.5\linewidth}{0.2pt} \\
    \vspace{0.2cm}
    {\small San Nicolás de los Garza, Nuevo León, México \\ Junio 2026}
\end{titlepage}

\newpage

% ==================== ÍNDICE ====================
\tableofcontents
\newpage

% ==================== SECCIONES ====================
\section*{Resumen}
\addcontentsline{toc}{section}{Resumen}

Este protocolo de investigación doctoral aborda el estudio teórico y computacional de la ingeniería de interfaces y la modulación de propiedades topológicas en heteroestructuras compuestas por aleaciones de Heusler y materiales bidimensionales tipo MXene. Mediante simulaciones de primeros principios basadas en la Teoría de la Funcional de la Densidad (DFT), se investigará la estabilidad estructural, la hibridación electrónica en la interfaz y la inducción de estados topológicos no triviales (como aislantes topológicos bidimensionales o semimetales de Weyl/Dirac). El objetivo principal es sintonizar estas propiedades mediante efectos de deformación (strain), dopaje y la aplicación de campos eléctricos externos, proporcionando así una base fundamental para el diseño de dispositivos espintrónicos y topológicos de próxima generación.

\vspace{0.5cm}
\noindent \textbf{Palabras clave:} Heteroestructuras, Aleaciones de Heusler, MXenes, Ingeniería de interfaces, Modulación topológica, DFT.

\vspace{1.5cm}

\section*{Abstract}
\addcontentsline{toc}{section}{Abstract}

This doctoral research protocol focuses on the theoretical and computational study of interface engineering and the modulation of topological properties in heterostructures composed of Heusler alloys and two-dimensional MXene materials. Using first-principles simulations based on Density Functional Theory (DFT), we will investigate structural stability, electronic hybridization at the interface, and the induction of non-trivial topological states (such as two-dimensional topological insulators or Weyl/Dirac semimetals). The primary goal is to tune these properties through strain engineering, doping, and the application of external electric fields, thereby providing a fundamental basis for the design of next-generation spintronic and topological devices.

\vspace{0.5cm}
\noindent \textbf{Keywords:} Heterostructures, Heusler alloys, MXenes, Interface engineering, Topological modulation, DFT.

\newpage

\section{Introducción y Antecedentes}
\label{sec:intro}

\subsection{Introducción}
En las últimas décadas, la física de la materia condensada ha experimentado una revolución con el descubrimiento de las fases topológicas de la materia, donde las propiedades globales de la estructura electrónica están protegidas contra perturbaciones locales y desorden por la simetría de inversión temporal u otras simetrías cristalinas~\cite{bansil2016colloquium}. 

Paralelamente, el estudio de heteroestructuras de baja dimensionalidad ha emergido como una de las plataformas más prometedoras para la ingeniería de materiales. Al combinar materiales con propiedades físicas distintas en una interfaz atómicamente plana, es posible inducir nuevos estados electrónicos y magnéticos que no están presentes en los materiales individuales. Este protocolo propone explorar la combinación de aleaciones de Heusler (conocidas por su alta polarización de espín y propiedades magnéticas controlables) y MXenes (carburos y nitruros de metales de transición bidimensionales con excelente conductividad y versatilidad superficial).

\subsection{Antecedentes}

\subsubsection{Aleaciones de Heusler}
Las aleaciones de Heusler son compuestos intermetálicos con estructura cúbica ($L2_1$ para Heusler completos o $C1_b$ para semi-Heusler) que exhiben una amplia variedad de propiedades físicas, desde ferromagnetismo de espín completo (half-metals) hasta superconductividad y comportamiento de aislantes topológicos~\cite{graf2011simple}. La posibilidad de sintonizar su brecha energética y la posición de sus niveles de energía mediante sustituciones químicas las convierte en candidatas ideales para acoplarse con materiales bidimensionales.

\subsubsection{MXenes}
Los MXenes son una familia de materiales bidimensionales descubiertos en 2011~\cite{naguib2011two}, con fórmula general $M_{n+1}X_n T_x$, donde $M$ representa un metal de transición temprana (como Ti, V, Cr), $X$ es carbono y/o nitrógeno, y $T_x$ representa grupos funcionales superficiales (como \ce{-O}, \ce{-F}, \ce{-OH}). Su gran área superficial, alta conductividad metálica e hidrofilia permiten una fácil integración en dispositivos nanoelectrónicos. Además, la presencia de metales de transición pesados en su estructura confiere un fuerte acoplamiento espín-órbita (SOC), el cual es un ingrediente clave para la inducción de estados topológicos no triviales.

\subsubsection{Heteroestructuras Heusler/MXene}
La interfase formada por una aleación de Heusler y un MXene ofrece un escenario sumamente rico para la modulación topológica. El fuerte acoplamiento de espín-órbita proveniente del MXene, combinado con el magnetismo y la polarización de espín de la aleación de Heusler, puede dar lugar a efectos físicos exóticos como el efecto Hall anómalo cuántico, la magnetorresistencia gigante o la generación de estados de borde topológicamente protegidos sintonizables. El control de la química y la estructura de la interfaz (ingeniería de interfaces) es, por lo tanto, el factor crucial para modular la topología de bandas en estas heteroestructuras.

\newpage

\section{Planteamiento del Problema, Hipótesis y Objetivos}
\label{sec:planteamiento}

\subsection{Planteamiento del Problema}
A pesar del enorme potencial de las heteroestructuras formadas por aleaciones de Heusler y materiales bidimensionales, la comprensión a nivel atómico y electrónico de sus interfaces sigue siendo limitada. En particular, la transferencia de carga en la interfaz, el desajuste de red (lattice mismatch), las reconstrucciones estructurales y la presencia de grupos funcionales en el MXene pueden alterar significativamente la estructura de bandas electrónicas, destruyendo o alterando los estados topológicos deseados. Por lo tanto, el problema científico radica en identificar cómo las características estructurales y químicas de la interfaz Heusler/MXene dictan las propiedades topológicas del sistema compuesto y determinar los mecanismos eficientes para su modulación y control externo.

\subsection{Pregunta de Investigación}
¿De qué manera influyen el ordenamiento atómico de la interfaz, el acoplamiento espín-órbita y los efectos de tensión epitaxial en la emergencia y modulación de fases topológicas no triviales en heteroestructuras compuestas de Heusler/MXene?

\subsection{Justificación}
El desarrollo de la espintrónica y la computación cuántica topológica requiere materiales con estados electrónicos altamente estables y protegidos contra la retrodispersión. Las heteroestructuras Heusler/MXene representan una nueva frontera para el diseño de materiales a medida. Comprender los principios físicos fundamentales que rigen estas interfaces mediante cálculos de primeros principios permitirá acelerar el diseño de dispositivos electrónicos sin disipación de calor, reduciendo el costo experimental y guiando la síntesis de nuevos materiales funcionales.

\subsection{Hipótesis}
La ingeniería de interfaces en heteroestructuras formadas por aleaciones de Heusler y MXenes modificados superficialmente permite inducir y sintonizar estados topológicos de borde no triviales mediante la transferencia selectiva de carga y el acoplamiento magnético interfacial, logrando una modulación dinámica a través de la aplicación de deformación uniaxial/biaxial y campos eléctricos externos.

\subsection{Objetivos}

\subsubsection{Objetivo General}
Estudiar teóricamente la estabilidad estructural, las propiedades electrónicas y la topología de bandas de las heteroestructuras Heusler/MXene mediante simulaciones de primeros principios basadas en DFT, con el fin de proponer mecanismos eficientes para la modulación y el control de fases topológicas.

\subsubsection{Objetivos Específicos}
\begin{enumerate}
    \item Diseñar y optimizar geométricamente modelos atómicos de interfaces Heusler/MXene, evaluando su estabilidad termodinámica y la energía de adhesión interfacial.
    \item Determinar la estructura de bandas electrónicas, la densidad de estados (DOS) y la polarización magnética local en la interfaz.
    \item Evaluar el efecto del acoplamiento espín-órbita (SOC) en la dispersión de bandas y calcular los invariantes topológicos (como el número de Chern o la Z2) para clasificar las fases electrónicas del sistema.
    \item Simular la influencia de estímulos externos (tensión mecánica epitaxial y campos eléctricos transversales) sobre la modulación de las propiedades topológicas y magnéticas de las heteroestructuras.
\end{enumerate}

\newpage

\section{Metodología y Detalles Computacionales}
\label{sec:metodologia}

La metodología propuesta para este proyecto doctoral se divide en tres etapas principales: diseño y construcción de los modelos de interfaces, optimización y cálculos DFT con Quantum ESPRESSO, y análisis de propiedades topológicas y de transporte.

\subsection{Construcción de Modelos de Heteroestructuras}
Se construirán superceldas que representen la interfaz entre la aleación de Heusler (plano de corte cristalográfico de interés, e.g., (001) o (111)) y el MXene bidimensional. Para minimizar los efectos de deformación artificial debido a la condición de frontera periódica en el plano, se seleccionarán combinaciones de superceldas tales que el desajuste de red sea inferior al 3\%:
\begin{equation}
    f = \frac{a_{\text{MXene}} - a_{\text{Heusler}}}{a_{\text{Heusler}}} \times 100\%
\end{equation}
Para evitar la interacción espuria entre imágenes periódicas en la dirección perpendicular a la interfaz, se introducirá una capa de vacío de al menos 18 \AA\ en la dirección $z$.

\subsection{Cálculos de Primeros Principios con Quantum ESPRESSO}
\label{subsec:qe}
Todas las simulaciones de estructura electrónica y optimización de geometría se realizarán utilizando el código de código abierto \textbf{Quantum ESPRESSO}~\cite{giannozzi2009quantum,giannozzi2017advanced}, el cual está basado en la Teoría de la Funcional de la Densidad (DFT), ondas planas y pseudopotenciales.

\subsubsection{Configuración de los Parámetros de Simulación}
\begin{itemize}
    \item \textbf{Funcional de Intercambio y Correlación:} Se empleará la aproximación de gradiente generalizado (GGA) con la parametrización de Perdew-Burke-Ernzerhof (PBE). En casos necesarios, se incorporarán correcciones de dispersión de van der Waals (e.g., DFT-D3) para modelar adecuadamente las interacciones débiles en la interfaz.
    \item \textbf{Pseudopotenciales:} Se utilizarán pseudopotenciales del tipo proyectores de ondas aumentadas (PAW) o ultrasuaves (USPP) completamente relativistas para describir la interacción del núcleo con los electrones de valencia, permitiendo capturar de manera precisa los efectos del acoplamiento espín-órbita (SOC).
    \item \textbf{Energía de Corte (Cutoff):} La energía de corte para la expansión de las funciones de onda (\texttt{ecutwfc}) y para la densidad de carga (\texttt{ecutrho}) se establecerá de acuerdo con las especificaciones de los pseudopotenciales elegidos para asegurar la convergencia total de la energía libre (típicamente \texttt{ecutwfc} $\ge$ 50 Ry y \texttt{ecutrho} $\ge$ 400 Ry).
    \item \textbf{Integración en la Zona de Brillouin:} Se utilizará una malla de puntos $k$ tipo Monkhorst-Pack para el muestreo de la zona de Brillouin, con una densidad adaptada al tamaño de la supercelda para asegurar la convergencia (por ejemplo, una malla de $8 \times 8 \times 1$ para optimizaciones de celda).
\end{itemize}

En la Tabla~\ref{tab:dft_parameters} se resumen los parámetros iniciales de simulación contemplados.

\begin{table}[htbp]
    \centering
    \caption{Parámetros iniciales de simulación para los cálculos DFT en Quantum ESPRESSO.}
    \label{tab:dft_parameters}
    \begin{tabular}{ll}
        \toprule
        \textbf{Parámetro} & \textbf{Configuración Propuesta} \\
        \midrule
        Software & Quantum ESPRESSO (v7.0+) \\
        Funcional XC & GGA-PBE (con correcciones vdW-DF / DFT-D3) \\
        Tipo de Pseudopotencial & PAW completamente relativista (con SOC) \\
        Corte de funciones de onda (\texttt{ecutwfc}) & 60 Ry \\
        Corte de densidad de carga (\texttt{ecutrho}) & 480 Ry \\
        Malla de puntos $k$ (Relajación) & $6 \times 6 \times 1$ (Monkhorst-Pack) \\
        Malla de puntos $k$ (Cálculos de DOS) & $12 \times 12 \times 1$ \\
        Criterio de convergencia de energía & $1.0 \times 10^{-6}$ Ry \\
        Criterio de convergencia de fuerzas & $1.0 \times 10^{-3}$ Ry/bohr \\
        \bottomrule
    \end{tabular}
\end{table}

\subsubsection{Cálculo de Propiedades Topológicas}
Para la clasificación topológica de las bandas, se utilizará la herramienta auxiliar \texttt{Wannier90} acoplada con Quantum ESPRESSO. A partir de las funciones de Wannier maximalmente localizadas (MLWFs), se construirán Hamiltonianos tight-binding efectivos para calcular los invariantes topológicos (como el número de Chern, la fase de Berry y los estados de borde en una geometría de cinta semi-infinita).

\newpage

\section{Cronograma de Actividades}
\label{sec:cronograma}

El proyecto está diseñado para desarrollarse en un periodo de 8 semestres (4 años), distribuidos de acuerdo con las etapas clave de investigación, análisis y difusión de resultados descritas en la Tabla~\ref{tab:cronograma}.

\begin{table}[htbp]
    \centering
    \caption{Cronograma semestral de actividades del proyecto de investigación doctoral.}
    \label{tab:cronograma}
    \resizebox{\textwidth}{!}{%
    \begin{tabular}{lcccccccc}
        \toprule
        \textbf{Actividad / Tarea} & \textbf{S1} & \textbf{S2} & \textbf{S3} & \textbf{S4} & \textbf{S5} & \textbf{S6} & \textbf{S7} & \textbf{S8} \\
        \midrule
        Revisión bibliográfica y actualización teórica & X & X & X & X & X & X & X & X \\
        Definición y modelado de interfaces Heusler/MXene & X & X & & & & & & \\
        Optimización geométrica y convergencia (DFT) & & X & X & & & & & \\
        Cálculos de estructura de bandas, DOS y magnetismo & & & X & X & & & & \\
        Cálculos relativistas (SOC) y análisis topológico & & & & X & X & X & & \\
        Simulación de modulación por strain y campos externos & & & & & X & X & & \\
        Escritura de manuscritos científicos (artículos) & & & & X & X & X & X & \\
        Redacción del documento de tesis y revisión & & & & & & & X & X \\
        Preparación y defensa del examen de grado & & & & & & & & X \\
        \bottomrule
    \end{tabular}%
    }
\end{table}

\noindent \textbf{Detalles de los hitos principales:}
\begin{itemize}
    \item \textbf{Fin del Año 1 (S2):} Modelos de interfaz listos y optimizados estructuralmente.
    \item \textbf{Fin del Año 2 (S4):} Propiedades magnéticas y electrónicas calculadas (envío de primer artículo).
    \item \textbf{Fin del Año 3 (S6):} Fases topológicas caracterizadas y mapeo del efecto de perturbaciones externas completado (envío de segundo artículo).
    \item \textbf{Fin del Año 4 (S8):} Redacción completa del borrador de tesis y examen de defensa doctoral.
\end{itemize}

\newpage

% ==================== REFERENCIAS ====================
\phantomsection
\addcontentsline{toc}{section}{Referencias Bibliográficas}
\printbibliography[title={Referencias Bibliográficas}]

\end{document}
